\section{辅助伴随表示“重夸克”方法}\label{sec:aux}

让我们从文献~\cite{Ji:2013dva} 定义的胶子准PDF出发：
\begin{align}\label{gluonpdfop}
\tilde g(x, P^z)=\int\frac{dz}{2\pi x n\cdot P}e^{i x z n\cdot P}\langle P|O(z,0)|P\rangle,\non\\
%\label{gluonpdfop}
O(z_2,z_1)=F_a^{z\mu}(z_2)L_{ab}(z_2, z_1)F_{b,\mu}^z(z_1).
\end{align}
$F_a^{\mu\nu}=\partial^\mu A_a^{\nu}-\partial^\nu A_a^\mu-g f_{abc}A_b^\mu A_c^\nu$ 是场强张量。除非另有说明，本节定义的量都是裸量。指标 $\mu$ 既可像胶子 PDF 中那样~\cite{Collins:1984xc} 只对横向分量求和，也可对所有洛伦兹分量求和。由于 $F^{zz}=0$，而 $F^{zt}$ 在大动量极限下受到幂次压低，两种求和方式都可用于定义胶子准PDF，且它们属于胶子准PDF算符的同一普适类（关于胶子极化算符普适类的讨论见文献~\cite{Hatta:2013gta}）。$P^\mu=(P^0, 0,0, P^z)$ 是外部强子的动量；$n^\mu=(0,0,0,1)$ 是指定规范连接 $L(z, 0) = {\cal P}\exp\left(-ig\int^z_0 d\lambda\,n^\mu  A_\mu(\lambda n)\right)$ 方向的类空矢量，其中 $A_\mu=A_{\mu, a}T_a$ 定义在伴随表示中。

另外，胶子准PDF也可以用基础表示的规范连接 $U(z_2,z_1)$ 来定义~\cite{Dorn:1981wa}
\beq\label{fundrep}
\hspace*{-1em}O(z_2,z_1)=2\,{\rm Tr}[F^{z\mu}(z_2)U(z_2,z_1)F_{z\mu}(z_1)U(z_1,z_2)],
\eeq
其中 $F^{\mu\nu}=F_{a}^{\mu\nu}t_a$，$t_a$ 是基础表示的生成元。
式~(\ref{gluonpdfop}) 便于研究重正化，而式~(\ref{fundrep}) 则更便于在格点上实现。下面我们主要聚焦于式~(\ref{gluonpdfop}) 的定义，但结果同样适用于式~(\ref{fundrep})。

为重正化 $O(z_2,z_1)$，我们仿照对夸克准PDF的研究~\cite{Ji:2017oey}，向 QCD 拉氏量引入一个记作 $\mathcal Q$ 的辅助伴随表示“重夸克”场。这个“重夸克”没有自旋自由度。含辅助场的拉氏量为
\beq\label{efflag}
\mathcal L=\mathcal L_{\rm{QCD}}+\overline{\mathcal Q}(x)(i n\cdot D-m)\mathcal Q(x),
\eeq
其中 $D_\mu=\partial_\mu+ig A_{\mu}$ 是伴随表示中的协变导数，并为辅助场 $\mathcal Q$ 引入了一个“剩余质量”，其理由稍后便会清楚。

借助辅助“重夸克” $\mathcal Q$，威尔逊线可以替换为两个此类场的乘积~\cite{Ji:2017oey}
\begin{align}\label{Eq:repl}
{\cal O}(z_2, z_1) &= F_a^{z\mu}(z_2){\mathcal Q}_a(z_2) \overline{\mathcal Q}_b(z_1) F^z_{b, \mu}(z_1)\non\\
&= J_1^{z\mu}(z_2) {\overline J}_{1, \mu}^z(z_1).
\end{align}
于是，式~(\ref{gluonpdfop}) 中非定域算符的重正化便归结为式~(\ref{Eq:repl}) 中两个定域复合算符（LCOs）乘积的重正化。
